\documentclass[10pt,respuestas,a4]{aleph-examen}
%\documentclass[11pt,a4]{aleph-examen}

% -- Paquetes adicionales 
\usepackage{aleph-comandos}
\usepackage{booktabs}
\usepackage{multicol}
\usepackage{enumitem}
\usepackage{array}

% -- Datos 
\institucion{Facultad de Ciencias Exactas, Naturales y Ambientales}
\carrera{Catálogo STEM}
\asignatura{Métodos Numéricos}
\tema{Examen}
\autor{Mario Cueva - Andrés Merino}
\fecha{Periodo 2026-1}

\logouno[0.18\textwidth]{Logos/logoPUCE_04_ac}
\definecolor{colortext}{HTML}{1957d1}
\definecolor{colordef}{HTML}{1957d1}
\fuente{montserrat}

\begin{document}

\encabezado

\section*{Indicaciones}
\begin{itemize}
    \item El examen tiene una duración de \textbf{90 minutos}.
    \item Se permite únicamente el uso de \textbf{calculadora científica}, lápiz y papel.
    \item Justifique todos los procedimientos. Las respuestas sin desarrollo tendrán puntaje parcial limitado.
    \item Redondee los resultados finales a \textbf{cuatro cifras decimales}, salvo que se indique lo contrario.
    \item En los métodos iterativos, presente las iteraciones en tablas ordenadas.
    \item En los ejercicios de integración y diferenciación, complete las tablas cuando sea necesario.
\end{itemize}

\section*{Examen}

\begin{preguntas}

%%%%%%%%%%%%%%%%%%%%%%%%%%%%%%%%%%%%%%%%
% Pregunta 1
%%%%%%%%%%%%%%%%%%%%%%%%%%%%%%%%%%%%%%%%

\item
Considere el siguiente sistema lineal: \puntaje{8}
\[
\begin{cases}
5x+y=12,\\
x+4y=10.
\end{cases}
\]

Aplique \textbf{tres iteraciones} de los métodos de \textbf{Jacobi} y \textbf{Gauss-Seidel}, partiendo de:
\[
(x^{(0)},y^{(0)})=(0,0).
\]

\begin{enumerate}[label=\alph*)]
    \item Despeje las ecuaciones para aplicar ambos métodos.
    \item Complete la tabla para el método de Jacobi.
    \item Complete la tabla para el método de Gauss-Seidel.
    \item Compare brevemente los resultados obtenidos.
\end{enumerate}

\[
x=\underline{\hspace{4cm}}, 
\qquad 
y=\underline{\hspace{4cm}}
\]

\begin{center}
\textbf{Método de Jacobi}

\renewcommand{\arraystretch}{1.4}
\begin{tabular}{c|c|c}
\toprule
$k$ & $x^{(k)}$ & $y^{(k)}$ \\
\midrule
0 & 0.0000 & 0.0000 \\
1 &  &  \\
2 &  &  \\
3 &  &  \\
\bottomrule
\end{tabular}
\end{center}

\begin{center}
\textbf{Método de Gauss-Seidel}

\renewcommand{\arraystretch}{1.4}
\begin{tabular}{c|c|c}
\toprule
$k$ & $x^{(k)}$ & $y^{(k)}$ \\
\midrule
0 & 0.0000 & 0.0000 \\
1 &  &  \\
2 &  &  \\
3 &  &  \\
\bottomrule
\end{tabular}
\end{center}

\vspace{0.3cm}

Comparación:

\[
\underline{\hspace{14cm}}
\]

\[
\underline{\hspace{14cm}}
\]

%%%%%%%%%%%%%%%%%%%%%%%%%%%%%%%%%%%%%%%%
% Pregunta 2
%%%%%%%%%%%%%%%%%%%%%%%%%%%%%%%%%%%%%%%%

\item
Construya el polinomio interpolante cuadrático de Newton para los puntos: \puntaje{6}
\[
(0,2), \qquad (1,5), \qquad (3,17).
\]

Luego estime el valor de la función en:
\[
x=2.
\]

Complete la tabla de diferencias divididas:

\begin{center}
\renewcommand{\arraystretch}{1.5}
\begin{tabular}{c|c|c|c}
\toprule
$i$ & $x_i$ & $f(x_i)$ & Diferencias divididas \\
\midrule
0 & 0 & 2 &  \\
1 & 1 & 5 &  \\
2 & 3 & 17 &  \\
\bottomrule
\end{tabular}
\end{center}

Escriba el polinomio de Newton:

\[
P_2(x)=\underline{\hspace{10cm}}
\]

Evalúe:

\[
P_2(2)=\underline{\hspace{5cm}}
\]

%%%%%%%%%%%%%%%%%%%%%%%%%%%%%%%%%%%%%%%%
% Pregunta 3
%%%%%%%%%%%%%%%%%%%%%%%%%%%%%%%%%%%%%%%%

\item
Ajuste una recta de regresión lineal de la forma: \puntaje{6}
\[
y=a_0+a_1x
\]
para los siguientes datos:

\[
\begin{array}{c|ccccc}
x_i & 1 & 2 & 3 & 4 & 5 \\
\hline
y_i & 3 & 4 & 6 & 8 & 9
\end{array}
\]

\begin{enumerate}[label=\alph*)]
    \item Complete la tabla de cálculo.
    \item Determine los coeficientes $a_0$ y $a_1$.
    \item Escriba el modelo ajustado.
    \item Estime el valor de $y$ cuando $x=6$.
\end{enumerate}

\begin{center}
\renewcommand{\arraystretch}{1.4}
\begin{tabular}{c|c|c|c}
\toprule
$x_i$ & $y_i$ & $x_i^2$ & $x_iy_i$ \\
\midrule
1 & 3 &  &  \\
2 & 4 &  &  \\
3 & 6 &  &  \\
4 & 8 &  &  \\
5 & 9 &  &  \\
\midrule
$\sum$ &  &  &  \\
\bottomrule
\end{tabular}
\end{center}

\[
a_1=\underline{\hspace{5cm}}, 
\qquad
a_0=\underline{\hspace{5cm}}
\]

Modelo:

\[
y=\underline{\hspace{7cm}}
\]

Estimación:

\[
y(6)=\underline{\hspace{5cm}}
\]

%%%%%%%%%%%%%%%%%%%%%%%%%%%%%%%%%%%%%%%%
% Pregunta 4
%%%%%%%%%%%%%%%%%%%%%%%%%%%%%%%%%%%%%%%%

\item
Para los datos: \puntaje{6}
\[
\begin{array}{c|cccc}
x_i & 0 & 1 & 2 & 3 \\
\hline
y_i & 2 & 3 & 6 & 11
\end{array}
\]

ajuste un modelo cuadrático de la forma:
\[
y=a_0+a_1x+a_2x^2
\]
utilizando la forma matricial de mínimos cuadrados.

\begin{enumerate}[label=\alph*)]
    \item Escriba la matriz de diseño $A$.
    \item Escriba el vector $\mathbf{y}$.
    \item Plantee el sistema normal $A^TA\mathbf{a}=A^T\mathbf{y}$.
    \item Determine los coeficientes del modelo.
\end{enumerate}

\[
A=
\begin{pmatrix}
 &  &  \\
 &  &  \\
 &  &  \\
 &  & 
\end{pmatrix},
\qquad
\mathbf{y}=
\begin{pmatrix}
\\
\\
\\
\end{pmatrix}
\]

\[
A^TA=
\begin{pmatrix}
 &  &  \\
 &  &  \\
 &  & 
\end{pmatrix},
\qquad
A^T\mathbf{y}=
\begin{pmatrix}
\\
\\
\end{pmatrix}
\]

Sistema normal:

\[
\begin{pmatrix}
 &  &  \\
 &  &  \\
 &  & 
\end{pmatrix}
\begin{pmatrix}
a_0\\
a_1\\
a_2
\end{pmatrix}
=
\begin{pmatrix}
\\
\\
\end{pmatrix}
\]

Modelo ajustado:

\[
y=\underline{\hspace{8cm}}
\]

%%%%%%%%%%%%%%%%%%%%%%%%%%%%%%%%%%%%%%%%
% Pregunta 5
%%%%%%%%%%%%%%%%%%%%%%%%%%%%%%%%%%%%%%%%

\item
Utilice diferenciación numérica para aproximar $f'(2)$ de la función: \puntaje{6}
\[
f(x)=\sqrt{x}
\]
con $h=0.1$, aplicando:

\begin{enumerate}[label=\alph*)]
    \item Diferencia hacia adelante.
    \item Diferencia hacia atrás.
    \item Diferencia centrada.
    \item Compare brevemente los resultados con el valor exacto:
    \[
    f'(2)=\frac{1}{2\sqrt{2}}.
    \]
\end{enumerate}

Complete la tabla:

\begin{center}
\renewcommand{\arraystretch}{1.4}
\begin{tabular}{c|c}
\toprule
$x$ & $f(x)=\sqrt{x}$ \\
\midrule
1.9 &  \\
2.0 &  \\
2.1 &  \\
\bottomrule
\end{tabular}
\end{center}

Diferencia hacia adelante:

\[
f'(2)\approx \underline{\hspace{6cm}}
\]

Diferencia hacia atrás:

\[
f'(2)\approx \underline{\hspace{6cm}}
\]

Diferencia centrada:

\[
f'(2)\approx \underline{\hspace{6cm}}
\]

Comparación:

\[
\underline{\hspace{14cm}}
\]

\[
\underline{\hspace{14cm}}
\]

%%%%%%%%%%%%%%%%%%%%%%%%%%%%%%%%%%%%%%%%
% Pregunta 6
%%%%%%%%%%%%%%%%%%%%%%%%%%%%%%%%%%%%%%%%

\item
Aproxime la integral: \puntaje{6}
\[
\int_0^2 (2x^2+1)\,dx
\]
mediante la \textbf{regla del trapecio compuesta} con $n=4$ subintervalos.

Luego calcule el error absoluto, considerando que el valor exacto de la integral es:
\[
\frac{22}{3}.
\]

\begin{enumerate}[label=\alph*)]
    \item Calcule el valor de $h$.
    \item Complete la tabla de valores.
    \item Aplique la fórmula del trapecio compuesto.
    \item Calcule el error absoluto.
\end{enumerate}

\[
h=\underline{\hspace{5cm}}
\]

\begin{center}
\renewcommand{\arraystretch}{1.4}
\begin{tabular}{c|c}
\toprule
$x_i$ & $f(x_i)=2x_i^2+1$ \\
\midrule
0.0 &  \\
0.5 &  \\
1.0 &  \\
1.5 &  \\
2.0 &  \\
\bottomrule
\end{tabular}
\end{center}

Aplicación de la fórmula:

\[
\int_0^2 (2x^2+1)\,dx
\approx
\underline{\hspace{8cm}}
\]

Error absoluto:

\[
E_a=\underline{\hspace{6cm}}
\]

%%%%%%%%%%%%%%%%%%%%%%%%%%%%%%%%%%%%%%%%
% Pregunta 7
%%%%%%%%%%%%%%%%%%%%%%%%%%%%%%%%%%%%%%%%

\item
Aproxime la integral: \puntaje{6}
\[
\int_0^2 (x^3+2x+1)\,dx
\]
mediante la \textbf{regla de Simpson $1/3$ simple}.

\begin{enumerate}[label=\alph*)]
    \item Identifique los valores de $a$, $b$ y el punto medio $m$.
    \item Evalúe la función en los tres puntos necesarios.
    \item Aplique la fórmula de Simpson $1/3$ simple.
    \item Indique si el resultado puede ser exacto para este tipo de función y explique brevemente.
\end{enumerate}

\[
a=\underline{\hspace{3cm}},
\qquad
b=\underline{\hspace{3cm}},
\qquad
m=\underline{\hspace{3cm}}
\]

\begin{center}
\renewcommand{\arraystretch}{1.4}
\begin{tabular}{c|c}
\toprule
$x$ & $f(x)=x^3+2x+1$ \\
\midrule
$a=$ &  \\
$m=$ &  \\
$b=$ &  \\
\bottomrule
\end{tabular}
\end{center}

Resultado:

\[
\int_0^2 (x^3+2x+1)\,dx
\approx
\underline{\hspace{7cm}}
\]

Explicación:

\[
\underline{\hspace{14cm}}
\]

\[
\underline{\hspace{14cm}}
\]

%%%%%%%%%%%%%%%%%%%%%%%%%%%%%%%%%%%%%%%%
% Pregunta 8
%%%%%%%%%%%%%%%%%%%%%%%%%%%%%%%%%%%%%%%%

\item
Aproxime la integral: \puntaje{6}
\[
\int_0^1 e^x\,dx
\]
mediante la \textbf{regla de Simpson $1/3$ compuesta} con $n=4$ subintervalos.

Además, explique por qué en este método el número de subintervalos debe ser par.

\begin{enumerate}[label=\alph*)]
    \item Calcule el valor de $h$.
    \item Complete la tabla de valores.
    \item Identifique los valores que se multiplican por $4$ y por $2$.
    \item Aplique la fórmula de Simpson $1/3$ compuesta.
    \item Explique por qué $n$ debe ser par.
\end{enumerate}

\[
h=\underline{\hspace{5cm}}
\]

\begin{center}
\renewcommand{\arraystretch}{1.4}
\begin{tabular}{c|c|c}
\toprule
$i$ & $x_i$ & $f(x_i)=e^{x_i}$ \\
\midrule
0 & 0.00 &  \\
1 & 0.25 &  \\
2 & 0.50 &  \\
3 & 0.75 &  \\
4 & 1.00 &  \\
\bottomrule
\end{tabular}
\end{center}

Valores con coeficiente $4$:

\[
\underline{\hspace{10cm}}
\]

Valores con coeficiente $2$:

\[
\underline{\hspace{10cm}}
\]

Resultado:

\[
\int_0^1 e^x\,dx
\approx
\underline{\hspace{7cm}}
\]

Explicación sobre $n$ par:

\[
\underline{\hspace{14cm}}
\]

\[
\underline{\hspace{14cm}}
\]

%%%%%%%%%%%%%%%%%%%%%%%%%%%%%%%%%%%%%%%%
% Pregunta 9
%%%%%%%%%%%%%%%%%%%%%%%%%%%%%%%%%%%%%%%%

\item
Pregunta conceptual de integración numérica: \puntaje{3}

Explique una diferencia entre la regla del trapecio compuesta y la regla de Simpson $1/3$ compuesta. Además, indique cuál de las dos suele dar una mejor aproximación cuando la función es suave y se usa el mismo número de subintervalos.

\[
\underline{\hspace{14cm}}
\]

\[
\underline{\hspace{14cm}}
\]

\[
\underline{\hspace{14cm}}
\]

%%%%%%%%%%%%%%%%%%%%%%%%%%%%%%%%%%%%%%%%
% Pregunta 10
%%%%%%%%%%%%%%%%%%%%%%%%%%%%%%%%%%%%%%%%

\item
Pregunta conceptual de métodos numéricos: \puntaje{3}

Explique por qué el método de Gauss-Seidel puede converger más rápido que el método de Jacobi en algunos sistemas lineales.

\[
\underline{\hspace{14cm}}
\]

\[
\underline{\hspace{14cm}}
\]

\[
\underline{\hspace{14cm}}
\]

\end{preguntas}

\end{document}

